\documentclass{beamer}

% You can also use a 16:9 aspect ratio:
%\documentclass[aspectratio=169]{beamer}
\usetheme{TACC16}
\usepackage{graphicx}

% It's possible to move the footer to the right:
%\usetheme[rightfooter]{TACC16}


% Detailed knowledge of application workload characteristics can
% optimize performance of current and future systems. This may sound
% daunting, with many HPC data centers hosting over 2,000 users running
% thousands of applications and millions of jobs per month.  XALT is an
% open source tool developed at the Texas Advanced Computing Center
% (TACC) that collects system usage information to quantitatively report
% how users are using your system. This session will explore the
% benefits of detailed application workload profiling and how the XALT
% tool has helped leading supercomputing sites unlock the power of their
% application usage data.

%% page
%\begin{frame}{}
%  \begin{itemize}
%    \item
%  \end{itemize}
%\end{frame}

%% page
%\begin{frame}[fragile]
%    \frametitle{}
% {\tiny
%    \begin{semiverbatim}
%    \end{semiverbatim}
%}
%  \begin{itemize}
%    \item
%  \end{itemize}
%
%\end{frame}



\begin{document}
\title[XALT]{Tracking imports for Python, R}
\author{Robert McLay}
\date{Feb. 17, 2022}

% page 1
\frame{\titlepage}

% page 2
\begin{frame}{XALT: Outline}
  \center{\includegraphics[width=.9\textwidth]{XALT_Sticker.png}}
  \begin{itemize}
    \item XALT can track executable that are run
    \item Also the shared libraries
    \item Can we track imports for Python and R?
    \item How could we do it?
  \end{itemize}
\end{frame}

% page 3
\begin{frame}{How can we track imported packages?}
  \begin{itemize}
    \item It will require special code unique for each tool
    \item Somehow we have to ``insert'' code into the import process
    \item This is typically accomplished by using some Hook provide by
      the tools' developers
  \end{itemize}
\end{frame}

% page 3
\begin{frame}{Tracking imports of packages for R}
  \begin{itemize}
    \item It started with R.
    \item James McComb \& Michael Scott from IU developed the R part
    \item They wrote code that intercepts the import action.
    \item XALT provides a program to take that data: xalt\_record\_pkg
    \item All packages tracker use this program to collect the data.
  \end{itemize}
\end{frame}

% page 4
\begin{frame}{XALT Prerequisites}
  \begin{itemize}
    \item A \texttt{path\_pattern} in a sites' Config.py.
    \item \texttt{ ['PKGS',  r'.*\/R'],}
    \item \texttt{ ['PKGS',  r'.*\/python[0-9.]*'],}
  \end{itemize}
\end{frame}

% page 5
\begin{frame}{XALT Env. Vars for PKGS}
  \begin{itemize}
    \item Since the execution is happenning during XALT Tracking
    \item The following environment variables are defined
      \begin{itemize}
        \item XALT\_DIR: The root directory of where xalt\_record\_pkg
        \item XALT\_RUN\_UUID: The run uuid for the current R or
          Python program running
      \end{itemize}
    \item The R and Python Import hooks only collect data if
      XALT\_RUN\_UUID is defined
  \end{itemize}
\end{frame}

% page 6
\begin{frame}{\texttt{xalt\_record\_pkg} usage}
  \begin{itemize}
    \item The hook routine does the following
    \item Gets XALT\_DIR and XALT\_RUN\_UUID from env.
    \item Builds command path to \texttt{xalt\_record\_pkg} using
      XALT\_DIR
    \item The rest of the command line is:
      \begin{itemize}
        \item -u $<$run\_uuid$>$
        \item program $<$name$>$
        \item xalt\_run\_uuid $<$run\_uuid$>$
        \item package\_name $<$pkg\_name$>$
        \item package\_version $<$pkg\_version$>$
        \item package\_path $<$pkg\_path$>$
      \end{itemize}
  \end{itemize}
\end{frame}

% page 7
\begin{frame}{\texttt{xalt\_record\_pkg} execution}
  \begin{itemize}
    \item \texttt{xalt\_record\_pkg} builds a json string w/ data
    \item Every import will generate a record
    \item Why?
  \end{itemize}
\end{frame}

% page 8
\begin{frame}{\texttt{xalt\_record\_pkg} execution (II)}
  \begin{itemize}
    \item The hook code in R, Python is called dynamically
    \item There is nothing recording that can be called at the end.
    \item Originally XALT was going to make PKGS not be sampled,
    \item Also generate a start record
    \item This way import records would have something to connnect
      with.
    \item However this is a bad idea!
  \end{itemize}
\end{frame}

% page 9
\begin{frame}{\texttt{xalt\_record\_pkg} execution (III)}
  \begin{itemize}
    \item XALT needs an execution record stored to save import data
    \item But there is too much Python runs to store every one
    \item Instead XALT uses /dev/shm in a unique directory (UUID)
    \item This avoids overlap with other executions
    \item But why?
  \end{itemize}
\end{frame}

% page 10
\begin{frame}{Why write package import data to /dev/shm?}
  \begin{itemize}
    \item Speed
    \item Python and R can be sampled
    \item Data is only sent on the ``wire'' at end of program if sampled
    \item Delete data otherwise.
  \end{itemize}
\end{frame}

% page 11
\begin{frame}{Sent on the ``wire''?}
  \begin{itemize}
    \item Import records are saved on /dev/shm
    \item This is independent of \$XALT\_TRANSMISSION\_STYLE
    \item At the end of the execution the end record and import
      records are sent via \$XALT\_TRANSMISSION\_STYLE
    \item This is only if sampled.
  \end{itemize}
\end{frame}

% page 12
\begin{frame}{What to do with this data?}
  \begin{itemize}
    \item Find the list of heavily imported packages
    \item Find who is using conda python
    \item XALT won't know if something is imported but not used
  \end{itemize}
\end{frame}


% page 20
\begin{frame}{Conclusions}
  \begin{itemize}
    \item We can extract useful data from XALT
    \item It is not quick on a VM.
    \item Would love to have the DB on SSD
    \item We use data to know what codes add to next benchmark.
  \end{itemize}
\end{frame}


% page 21
\begin{frame}{Future Topics?}
  \begin{itemize}
    \item Others?
  \end{itemize}
\end{frame}
%

%
\begin{frame}{License}
\footnotesize
\textcopyright The University of Texas at Austin, \the\year\\
\vspace{0.5 cm}
This work is licensed under the Creative Commons Attribution Non-Commercial 3.0 Unported License. To view a copy of this license, 
visit http://creativecommons.org/licenses/by-nc/3.0/\\
\vspace{0.5 cm}
When attributing this work, please use the following text: 
"Designing and Administering Large-scale Systems", Texas Advanced Computing Center, \the\year. Available under a Creative Commons Attribution Non-Commercial 3.0 Unported License.
\newline
\newline
\includegraphics[width=1cm]{./figures/cc.png}
\end{frame}

\end{document}
